%% Abstract body only -- no heading, no size command. Included by paper.tex.
The LIAR benchmark includes five ``credit history'' columns per speaker, but the
publication describing the data set notes that these include the statement to be
predicted, thus emphasizing that these should only be used after the current
annotation has been subtracted from it. We provide the evidence for the cost of
following these instructions. Forests fed with only those five columns from the
data set achieve 0.394 accuracy at the test time in the six-way classification
task, but the amounts achieved in the experiment when highlighting the value of
those columns according to what was stated in the paper when subtracting the
current label are equal to 0.659 against the majority class accuracy of 0.208.
Our proof comes from the counting rather than reasoning related to the algorithm
as to how it works. There are 3\,318 speakers who share 392 distinct counting
vectors in total, meaning that the vector used for this subtraction can be
treated as a deviation from a landmark that can be remembered by the classifier
and that this vector indicates the actual label as it retrieves the same value
for 36.2\% of the test rows. Subtraction is insufficient for the independent
reasons since the counts in question may be correlated as they are aggregated
across training, validation and test samples for 2\,846 of the defined number of
speakers. We then reconstruct speaker history according to three principles as
far as the training labels, out-of-fold coding and prior smoothing are
concerned, thus demonstrating that the model returns 0.228 score according to
the history approach, which is better than the primary class score. Finally, we
test six classifiers in a set of experiments under a protocol that selects on
validation and scores on testing once; the best configurations in this case
return 0.285 for six-class and 0.673 for binary accuracy, and the hybrid
configuration has performed better than both single modalities in cases.
